\documentclass[11pt]{article}

\usepackage[letterpaper,margin=0.9in]{geometry}
\usepackage{amsmath,amssymb,amsfonts}
\usepackage{array,booktabs,longtable,tabularx}
\usepackage{enumitem}
\usepackage{titlesec}
\usepackage{fancyhdr}
\usepackage{lastpage}
\usepackage{hyperref}
\usepackage{xcolor}
\usepackage{graphicx}

\setlength{\headheight}{14pt}
\hypersetup{
  colorlinks=true,
  linkcolor=black,
  urlcolor=black,
  citecolor=black
}

\pagestyle{fancy}
\fancyhf{}
\lhead{White Sphere--Cylinder Casimir Probe}
\rhead{Source-Function Audit}
\cfoot{\thepage\ of \pageref{LastPage}}

\titleformat{\section}{\large\bfseries\uppercase}{\thesection.}{0.5em}{}
\titleformat{\subsection}{\normalsize\bfseries}{\thesubsection}{0.5em}{}
\setlist[itemize]{nosep,leftmargin=1.2em}
\setlist[enumerate]{nosep,leftmargin=1.4em}
\renewcommand{\arraystretch}{1.16}

\newcommand{\um}{\,\mu\mathrm{m}}
\newcommand{\R}{\mathbb{R}}
\newcommand{\shell}{\mathcal{S}_{\mathrm{shell}}}
\newcommand{\pair}{\mathcal{P}_{\mathrm{sw}}}
\newcommand{\Eproxy}{\mathcal{E}_{\mathrm{proxy}}}
\newcommand{\Nneg}{M_{-}}

\title{\bfseries A Role-Resolved Source-Function Audit of the White et al. Sphere--Cylinder Casimir Geometry}
\author{
Daniel S. Goldman\\
Independent researcher\\
\href{https://orcid.org/0000-0003-2835-3521}{ORCID: 0000-0003-2835-3521}
}
\date{\today}

\begin{document}
\maketitle

\begin{abstract}
White et al.'s sphere-in-cylinder Casimir result contains two separable parts:
a boundary-geometry observation and a warp-source interpretation. The
boundary-geometry question concerns shape: whether a $1\um$ diameter sphere
centered in a nominal $4\um$ diameter cylinder organizes the vacuum-boundary
proxy into the annular region where the reported warp-like shell appears. The
source-function question concerns role and scale: whether that organized
region remains a shell-placed source proxy under channel accounting, how its
calibrated magnitude compares with an Alcubierre source demand, and how the
central readout path participates.

This paper evaluates the claim ladder with three role-resolved proxy layers. A
pair-resolved v-loop scalar morphology probe supplies the shape read: each
dimensionless contribution is a loop-center/scale event whose segments cross
both the sphere surface and the inner cylinder-wall surface, followed by the
declared scale weight. A finite-difference tensor-scale audit supplies
the source-function read: common random loops, parallel-plate Casimir
calibration brackets, source-demand comparison, linearized metric bounds,
material-energy estimates, and explicit readout accounting. A synthetic
readout-discrimination layer then treats the finite-difference shell response
as a geometry-coded signal template and tests recovery against ordinary
electromagnetic, material, drift, vibration, and readout nuisance templates.
The sphere, cylinder wall, annular shell, central readout path, boundary
infrastructure, ordinary electromagnetic competitor region, and far-field
control are scored as separate roles.

The integrated result separates a positive morphology/source-placement result
from a closed direct-readout interpretation under the modeled laboratory
backgrounds. The selective scalar channel concentrates signed proxy magnitude
in the annular source region, places the radial maximum near the sphere-wall
gap, and keeps far-field leakage small. The tensor-scale proxy preserves this
shell placement and assigns a small fraction of the source proxy to the central
readout role. This supports the shell-morphology interpretation and clarifies
the readout path as a measurement role distinct from the main source channel.
SI calibration then fixes the magnitude layer: the shell-shaped source proxy
remains at Casimir gravitational strength, far beneath representative
warp-source demand and far beneath a recoverable timing/readout scale for the
modeled apparatus. The synthetic discrimination layer turns that scale result
into a laboratory gate. Geometry-coded perturbations can separate a shell
template from ordinary electromagnetic and material nuisance templates once a
shell-coded observable has already reached background-scale amplitude.
White's geometry therefore reads as a real shell-shaped Casimir boundary
morphology with useful source-placement structure; the current direct
laboratory detection interpretation closes under this audit class.
\end{abstract}

\section{Question}

White et al. applied worldline numerics to custom Casimir geometries and found
a negative energy-density morphology that they described as intersecting with
the energy-density requirements of the Alcubierre metric \cite{White2021}. The
most compact version of that reported geometry is the sphere-in-cylinder toy model: a
$1\um$ diameter sphere centered in a $4\um$ diameter cylinder. The reported
calculation produced a shell-like negative distribution that was visually close
to a two-dimensional section of an Alcubierre-style negative-energy shell.

This study evaluates the claim ladder behind that comparison. The first layer
is scalar morphology: whether the sphere-cylinder boundary geometry organizes
signed pair occupancy into the proposed shell under scale, seed, wall, and
geometry variations. Here ``robust'' means concentration in the intended
sphere-cylinder shell region under baseline seed changes, wall-thickness
changes, reported scale windows, and single-seed modest sphere and cylinder
perturbations.

The second layer is source-function relevance. A scalar shell silhouette and
an Alcubierre source occupy different claim layers. The stronger question asks
where the role-accounted morphology sits under finite-difference source
proxies, how the calibrated magnitude compares with a representative
Alcubierre source demand, how the central readout path participates, and how
material and ordinary electromagnetic competitors enter a laboratory
interpretation. The evidence is therefore organized as morphology, role
placement, scale closure, material competition, and readout accounting.

The third layer is readout discrimination. A geometry-coded laboratory schedule
can vary the same apparatus controls that move the shell proxy and the ordinary
electromagnetic competitors. The discriminating question asks what
signal-to-background ratio allows a shell-template coefficient to be recovered
from synthetic readout data with low ordinary-EM-only false positives. This
layer assigns the perturbation schedule an engineering role: channel separation
across controlled geometry changes.

The analysis protocol draws on prior source-placement and role-separation work
\cite{GoldmanRail2026,GoldmanDisclosure}. The White device is treated here in
its own Casimir terms: roles, channels, controls, and observables are separated
before source or transport language is assigned.

\section{Geometry and Roles}

The evaluated geometry is the White et al. sphere-cylinder toy model:
\begin{equation}
  d_s = 1.0\um,\qquad d_c = 4.0\um,
\end{equation}
where $d_s$ and $d_c$ are nominal diameters and the sphere is centered in the
cylinder. The crossed cylinder surface is the inner wall surface
\begin{equation}
  R_{\mathrm{inner}} = R_c - \frac{t_{\mathrm{wall}}}{2}.
\end{equation}
For the $d_c=4.0\um$ baseline, the reported $0.2\um$ and $0.4\um$ wall rows
use inner crossing diameters of $3.8\um$ and $3.6\um$. The analysis
treats the cylinder wall as boundary infrastructure, the sphere as the
stress-shaping body, the annular region outside the sphere and inside the
cylinder as the source-shell candidate, the central axis as the proposed
transit/readout channel, and the exterior region as far-field control.

For a grid point $x=(x,y,z)$ in a two-dimensional section, let $r_s(x)$ be the
distance from the sphere center and let $r_\perp(x)=\sqrt{y^2+z^2}$ be the
transverse radius relative to the cylinder axis. The near-shell band used for
the principal concentration metric is
\begin{equation}
  \shell =
  \{x:\ R_s < r_s(x) \leq R_s + 0.75\um,\quad
  r_\perp(x) < R_c - 0.10\um\}.
\end{equation}
This fixed nominal band is a predefined role-labeled accounting region used for
all reported cases to represent the intended sphere-cylinder shell. The
material wall, central readout mask, and far field are tracked with their own
accounting masks, and those masks may overlap on grid cells where geometric
roles share space.

\section{Loop Construction and Kernel}

The loop generator follows the v-loop construction quoted in White et al. The
sampled Gaussian variables use the paper's $\exp(-w_i^2)$ convention, i.e. a
normal distribution with variance $1/2$. Let $y_{a,n}\in\R^3$ denote point
$n$ on loop $a$, with the loop center of mass subtracted.

The scalar proxy uses segment/surface crossings. For each translated and
scaled loop
\begin{equation}
  \Gamma_{a,\lambda,c}=\{c+\lambda y_{a,n}\}_{n=1}^{N},
\end{equation}
the kernel determines whether any segment of the loop crosses the sphere surface
and whether any segment crosses the inner cylinder wall. A loop contributes to
the pair interaction when both crossings occur:
\begin{equation}
  I_{a,\lambda}(c)=
  \mathbf{1}\{\Gamma_{a,\lambda,c}\cap \partial B_s \neq \varnothing\}
  \mathbf{1}\{\Gamma_{a,\lambda,c}\cap \partial C_{\mathrm{inner}}\neq \varnothing\}.
\end{equation}
The normalized scalar morphology proxy at grid center $c$ is then
\begin{equation}
  \Eproxy(c) =
  -\frac{1}{N_{\mathrm{loop}}}
  \sum_{\lambda\in\Lambda}
  \frac{1}{\lambda^4}
  \sum_{a=1}^{N_{\mathrm{loop}}} I_{a,\lambda}(c).
\end{equation}
The $\lambda^{-4}$ weighting keeps the scalar proxy aligned with the expected
scale sensitivity of a four-dimensional Casimir-type contribution. The sum over
$\lambda$ is a declared discrete geometry-scale score for comparing morphology
within the same support interval. A stress-tensor or energy-density estimator
adds continuum proper-time quadrature, quadrature weights, point-splitting or
derivative operators, material/boundary response, and SI normalization. The
output category here is a normalized scalar morphology proxy with explicit
dimensionless occupancy status.

The single-body controls are embedded in the scoring rule. Contributions are
accepted when a loop crosses both the sphere and the inner cylinder surface.
This makes the probe a direct pair-interaction read; generic
boundary contacts carry zero contribution in this score.

\section{Readout Metrics}

Let
\begin{equation}
  \Nneg(c)=\max[-\Eproxy(c),0]
\end{equation}
be the signed-occupancy magnitude recorded by the proxy. The principal shell
concentration metric is
\begin{equation}
  F_{\shell} =
  \frac{\sum_{c\in\shell}\Nneg(c)}
       {\sum_{c}\Nneg(c)}.
\end{equation}
The shell enrichment is the ratio between mean signed magnitude in the
near-shell band and mean signed magnitude outside the band:
\begin{equation}
  Q_{\shell} =
  \frac{\langle \Nneg\rangle_{\shell}}
       {\langle \Nneg\rangle_{\R^2\setminus\shell}}.
\end{equation}
The top-decile shell fraction records the fraction of the strongest signed
pixels that fall in $\shell$. Boundary, far-field, and readout participation
are computed as the corresponding magnitude fractions in their role-labeled
regions. Two central-channel quantities are tracked: the explicit readout path
from the geometry role label, with radius $0.05\um$, and a wider central-axis
mask,
\begin{equation}
  A_{0.15}=\{x:\ r_\perp(x)\leq 0.15\um\}.
\end{equation}
The explicit readout path is an accounting mask applied after the scalar field
is computed. The crossing geometry keeps the sphere and cylinder boundary model
fixed; a modeled bore, conductor, material removal, or propagation observable is
a separate readout calculation. The wider $A_{0.15}$ mask records nearby
central-axis participation on the grid.

\section{Computation Configuration}

The calculation used the configuration in Table~\ref{tab:config}. The
computational design emphasizes the interaction kernel and role-accounted scale
windows. Loop count and grid size set the scope of the present scalar
morphology calculation. A normalized White-style scalar comparison adds a larger
loop ensemble, continuum proper-time scaling, parallel-plate normalization, and
direct comparison to the reported $2000$-loop run. The source-function layer
below adds a local finite-difference tensor-scale proxy, SI calibration
brackets, source-demand comparison, material estimates, and readout accounting.
The synthetic discrimination layer adds template recovery under geometry-coded
ordinary electromagnetic and material nuisance backgrounds.
A full worldline stress tensor and a modeled laboratory propagation observable
belong to a higher-fidelity physical layer.
The validation scope for this manuscript is crossing geometry, sign accounting,
role bookkeeping, scale-window morphology, source-role stability, and calibrated
scale comparison, with synthetic readout discrimination reported as an
identifiability threshold.

\begin{table}[h]
\centering
\caption{Pair-surface scalar-probe configuration.}
\label{tab:config}
\begin{tabularx}{\linewidth}{>{\raggedright\arraybackslash}p{0.34\linewidth}X}
\toprule
Item & Value \\
\midrule
Kernel & Pair-resolved segment/surface crossing \\
Surfaces & Sphere surface and inner cylinder wall \\
Loop method & Paper-style v-loop \\
Baseline seeds & $1,7,17$ \\
Wall thicknesses & $0.2\um$, $0.4\um$ \\
Inner crossing diameters for $d_c=4.0\um$ & $3.8\um$, $3.6\um$ \\
Baseline scale windows & $0.75$--$1.50\um$, $1.50$--$2.75\um$, $0.25$--$5.00\um$ \\
Geometry perturbations & sphere $0.9\um$, sphere $1.1\um$, cylinder $3.8\um$, cylinder $4.2\um$ \\
Perturbation seed & $1$ \\
Loops & $96$ \\
Points per loop & $192$ \\
Grid & $37\times37$ over $\pm2.5\um$ \\
Workers & $4$ \\
\bottomrule
\end{tabularx}
\end{table}

\section{Baseline Results}

Table~\ref{tab:baseline} gives the baseline aggregate over the three seeds.
The mid-scale window is the most spatially selective morphology channel. It covers about
$0.35$--$0.37$ of the sampled pixels with signed occupancy values, yet it places
$0.647$--$0.726$ of the signed magnitude in the fixed nominal near-shell band.
Boundary and far-field fractions are negligible, and the peak radial bin
remains at $1.125\um$.

\begin{table}[h]
\centering
\caption{Baseline pair-surface scalar-probe read. Values are means over seeds $1,7,17$.}
\label{tab:baseline}
\small
\begin{tabular}{lrrrrrrr}
\toprule
Window & Wall & Neg. frac. & $F_{\shell}$ & $Q_{\shell}$ & Top-decile shell & Boundary & Far \\
\midrule
$0.75$--$1.50$ & $0.2$ & $0.347699$ & $0.646612$ & $9.782$ & $0.781228$ & $0.000283$ & $0.000000$ \\
$0.75$--$1.50$ & $0.4$ & $0.365961$ & $0.726420$ & $14.206$ & $0.883074$ & $0.001319$ & $0.000000$ \\
$1.50$--$2.75$ & $0.2$ & $0.892866$ & $0.540825$ & $6.289$ & $0.771857$ & $0.012854$ & $0.005076$ \\
$1.50$--$2.75$ & $0.4$ & $0.893109$ & $0.567666$ & $7.012$ & $0.875146$ & $0.023294$ & $0.000718$ \\
$0.25$--$5.00$ & $0.2$ & $1.000000$ & $0.504602$ & $5.437$ & $0.708029$ & $0.020623$ & $0.014162$ \\
$0.25$--$5.00$ & $0.4$ & $1.000000$ & $0.547679$ & $6.464$ & $0.790754$ & $0.032290$ & $0.002884$ \\
\bottomrule
\end{tabular}
\end{table}

The upper-scale window is broader and fills most of the section, and its
largest signed pixels remain shell-concentrated. The broad scale-sum window
gives apparatus-wide signed scalar participation because larger loops connect
the sphere-wall system across much of the sampled section. Even in the broad
field, approximately half of the signed magnitude remains in the fixed nominal
near-shell band with $5$--$6$ fold enrichment. These rows separate two scalar
channels: apparatus-wide broad-window participation and localized mid-scale
pair interaction.

Table~\ref{tab:readout} gives the corresponding central-channel read. In the
mid-scale channel, the explicit readout path carries $0.000651$--$0.001225$ of
the signed proxy magnitude, and the wider central-axis mask carries
$0.002918$--$0.005459$. The upper-scale and broad channels add apparatus-wide
central participation at the few-percent level on the explicit readout path.

\begin{table}[h]
\centering
\caption{Baseline readout-axis participation. Values are means over seeds $1,7,17$.}
\label{tab:readout}
\small
\begin{tabular}{lrrr}
\toprule
Window & Wall & Explicit readout path & Central axis $A_{0.15}$ \\
\midrule
$0.75$--$1.50$ & $0.2$ & $0.000651$ & $0.002918$ \\
$0.75$--$1.50$ & $0.4$ & $0.001225$ & $0.005459$ \\
$1.50$--$2.75$ & $0.2$ & $0.019870$ & $0.063433$ \\
$1.50$--$2.75$ & $0.4$ & $0.023333$ & $0.073679$ \\
$0.25$--$5.00$ & $0.2$ & $0.023119$ & $0.071958$ \\
$0.25$--$5.00$ & $0.4$ & $0.023688$ & $0.073302$ \\
\bottomrule
\end{tabular}
\end{table}

\section{Geometry Perturbations}

Table~\ref{tab:perturb} summarizes the mid-scale perturbation cases at seed
$1$. In this single-seed geometry-direction check, the tighter cylinder
strengthens the near-shell concentration, the wider cylinder weakens it, and
both sphere-size perturbations retain a recognizable shell-localized response.
The peak radial bin remains near $1.125\um$ in the mid-scale cases.

\begin{table}[h]
\centering
\caption{Mid-scale perturbation read at seed $1$.}
\label{tab:perturb}
\small
\begin{tabular}{lrrrrr}
\toprule
Geometry & Wall & Neg. frac. & $F_{\shell}$ & $Q_{\shell}$ & Top-decile shell \\
\midrule
Cylinder $3.8\um$ & $0.2$ & $0.379839$ & $0.718409$ & $13.618$ & $0.865385$ \\
Cylinder $3.8\um$ & $0.4$ & $0.385683$ & $0.782528$ & $19.208$ & $0.962264$ \\
Cylinder $4.2\um$ & $0.2$ & $0.310446$ & $0.539340$ & $6.250$ & $0.627907$ \\
Cylinder $4.2\um$ & $0.4$ & $0.372535$ & $0.632558$ & $9.189$ & $0.745098$ \\
Sphere $0.9\um$ & $0.2$ & $0.344777$ & $0.628413$ & $9.658$ & $0.708333$ \\
Sphere $0.9\um$ & $0.4$ & $0.352082$ & $0.719429$ & $14.643$ & $0.897959$ \\
Sphere $1.1\um$ & $0.2$ & $0.395179$ & $0.689049$ & $10.860$ & $0.818182$ \\
Sphere $1.1\um$ & $0.4$ & $0.401753$ & $0.762371$ & $15.723$ & $0.909091$ \\
\bottomrule
\end{tabular}
\end{table}

The upper-scale perturbation cases retain a broader and still organized read:
near-shell magnitude fractions range from $0.515$ to $0.595$, top-decile
near-shell fractions range from $0.719$ to $0.942$, and far-field magnitude
fractions stay below $0.006655$. The perturbation rows give the same
geometry-direction picture as the baseline calculation: the mid-scale channel
is localized, and larger scale channels add apparatus-wide participation while
the largest magnitudes remain near the shell. Across the perturbation cases, the mid-scale explicit
readout-path fraction stays in the $0.000000$--$0.002634$ range, with the wider
central-axis fraction in the $0.000000$--$0.008408$ range. The upper-scale
perturbation cases carry $0.017245$--$0.027398$ on the explicit readout path
and $0.055520$--$0.084184$ on the wider central-axis mask.

\section{Tensor-Scale Source-Function Probe}

The source-function layer keeps the same role separation and asks a stronger
question: whether the scalar morphology remains source-shell placed under
paired finite-difference source proxies, whether a calibrated magnitude is
competitive with an Alcubierre source demand, and whether the central readout
path carries the source proxy. The calculation is a finite-difference
source-channel proxy. Its role is to bind the morphology result to transparent
scale and channel accounting before a full worldline stress tensor is attempted.

The focused tensor-scale calculation used a $73\times73$ grid over $\pm2.5\um$, $12$
loop blocks, $32$ loops per block, $256$ points per loop, the two selective
scale windows $0.75$--$1.50\um$ and $1.50$--$2.75\um$, wall thicknesses
$0.2\um$ and $0.4\um$, a $0.025\um$ finite-difference step, and the same
$0.05\um$ readout accounting radius. It produced $48$ tensor blocks, $4032$
stress-channel rows, $1536$ source-demand rows, and $576$ material-competitor
rows. The calibration stage used parallel-plate Casimir brackets for ideal EM
and scalar Dirichlet families over gaps $0.75$--$2.75\um$; gap ordering passed
in all calibration families and scale windows.

Table~\ref{tab:tensor_roles} gives the baseline role fractions for the
negative-magnitude source proxy. The mid-scale channel is almost entirely a
source-shell channel. The upper-scale channel is broader, with larger
stress-body and ordinary electromagnetic competitor participation, and it
remains shell dominated.

\begin{table}[h]
\centering
\caption{Tensor-scale role fractions for the finite-difference source proxy.}
\label{tab:tensor_roles}
\small
\begin{tabular}{lrrrrrr}
\toprule
Window & Wall & Shell & Readout & EM region & Body & Far \\
\midrule
$0.75$--$1.50$ & $0.2$ & $0.995759$ & $0.000231$ & $0.000679$ & $0.004074$ & $0.000000$ \\
$0.75$--$1.50$ & $0.4$ & $0.992679$ & $0.000340$ & $0.001962$ & $0.007203$ & $0.000000$ \\
$1.50$--$2.75$ & $0.2$ & $0.938713$ & $0.009849$ & $0.019400$ & $0.050144$ & $0.003941$ \\
$1.50$--$2.75$ & $0.4$ & $0.932350$ & $0.011070$ & $0.031461$ & $0.057991$ & $0.000807$ \\
\bottomrule
\end{tabular}
\end{table}

The finite-difference channel read identifies the geometry controls that move
the proxy. The source-shell channel responds primarily to cylinder inner radius
and sphere diameter. In the mid-scale rows, the cylinder-radius shell medians
are $541.8$ and $739.5$ per micrometer for wall thicknesses $0.2\um$ and
$0.4\um$, while the sphere-diameter shell medians are $-345.0$ and $-492.1$
per micrometer. In the upper-scale rows, the corresponding shell medians are
$366.5$, $386.9$, $-314.1$, and $-340.6$ per micrometer. Cylinder length is
essentially null at this scale, and perturbing the readout accounting radius
acts on the readout channel without generating a source-shell response. That
pattern anchors the shell result to the sphere-cylinder boundary geometry and
keeps readout-mask bookkeeping in its own accounting role.

Table~\ref{tab:scale_bound} summarizes the calibrated scale comparison. The
best source-demand ratio occurs in the ideal-EM mid-scale family at material
factor $1.0$, target width $2.0\um$, and speed parameter $10^{-9}$, and is
$1.019\times10^{-39}$. The largest linearized timing bound is
$7.81\times10^{-76}\,\mathrm{s}$. These numbers are scale bounds. A transport
prediction requires a propagation model, while the present source-demand table
places the shell-like morphology many orders below Alcubierre source-demand
magnitude in this harness.

\begin{table}[h]
\centering
\caption{Calibrated source and metric-scale bounds.}
\label{tab:scale_bound}
\small
\begin{tabular}{lr}
\toprule
Quantity & Value \\
\midrule
Maximum source-demand ratio & $1.019\times10^{-39}$ \\
Full source-demand ratio range & $3.20\times10^{-62}$--$1.019\times10^{-39}$ \\
Integrated shell energy range & $5.24\times10^{-25}$--$2.35\times10^{-21}\,\mathrm{J}$ \\
Readout-coupled energy range & $0$--$1.42\times10^{-23}\,\mathrm{J}$ \\
Maximum gravitational radius & $2.34\times10^{-67}\,\mathrm{m}$ \\
Maximum fractional path scale & $2.93\times10^{-62}$ \\
Maximum timing bound & $7.81\times10^{-76}\,\mathrm{s}$ \\
\bottomrule
\end{tabular}
\end{table}

The readout summary distinguishes source-placement accounting from a later
laboratory propagation measurement. Source placement asks how much of the proxy
is assigned to the central path; propagation asks how photons, electrons, or
current respond to the surrounding material geometry. This source-function
probe directly tests source placement. Propagation observables require their
own optical, electronic, or current-transport calculation. In the
calibrated summary, the mid-scale readout source fraction is zero; in the
upper-scale summary, the median readout source fraction is $0.010449$ while the
shell channel fraction is $0.934467$, a shell/readout ratio of about $89.4$.
The source proxy is therefore shell-first, with readout retained as a distinct
measurement role.

The material and ordinary electromagnetic estimates set the laboratory
interpretation boundary. Across the material table, the Casimir-to-material
rest-energy ratio ranges from $1.06\times10^{-29}$ to $3.95\times10^{-25}$,
and the patch-potential-to-Casimir energy ratio ranges from $2.57$ to
$1.15\times10^{6}$. Patch-potential and contact-loading flags appear in every
material row. These estimates leave the shell morphology intact and set a
strong condition on experimental interpretation: ordinary material and
electromagnetic effects require explicit modeling.

\section{Synthetic Readout Discrimination}

The Stage 3 readout and material results motivate a direct engineering
question: can a geometry-modulation schedule distinguish a shell-channel
response from ordinary laboratory backgrounds? The synthetic discrimination
layer treats the Stage 3 finite-difference shell response as a template and
mixes it with nuisance templates for capacitance shift, patch potential,
contact loading, material response, thermal drift, vibration alignment,
waveguide dispersion, and readout-circuit artifact. The output is an
identifiability threshold for a shell-coded observable under a declared
geometry schedule.

The upgraded synthetic run uses $48$ block-level shell-template samples from
\path{tensor/stress_channel_points.parquet} and a geometry-derived nuisance
model. The fit template is the median Stage 3 shell response. Each synthetic
dataset draws a shell template from the block-level sample set. This tests
recovery under template mismatch as well as ordinary nuisance mixing. The
median correlation between sampled shell templates and the median fit template
is $0.884$, and the largest shell--nuisance template correlation is $0.928$.
The latter value occurs because geometry-derived patch, capacitance, material,
and readout effects share several apparatus controls with the shell template.

\begin{table}[h]
\centering
\caption{Synthetic readout-discrimination recovery under block-bootstrap shell templates and geometry-derived nuisance templates.}
\label{tab:synthetic_discrimination}
\small
\begin{tabular}{rrrr}
\toprule
Signal/background & Shell recovery & EM-only false positive & Median shell $z$ \\
\midrule
$1.0$ & $0.705$ & $0.040$ & $2.801$ \\
$1.2$ & $0.845$ & $0.040$ & $3.178$ \\
$1.4$ & $0.875$ & $0.015$ & $3.933$ \\
$1.6$ & $0.915$ & $0.045$ & $4.182$ \\
$1.8$ & $0.950$ & $0.030$ & $5.002$ \\
$2.0$ & $0.975$ & $0.035$ & $5.305$ \\
\bottomrule
\end{tabular}
\end{table}

The discrimination threshold is therefore near signal-to-background ratio
$1.2$ for the current synthetic gate. The full geometry schedule is the
preferred schedule in this model; reduced schedules lose rank and carry larger
shell--nuisance degeneracies. Geometry modulation supplies separation by
forcing shell, material, and readout templates to move through different
patterns across sphere size, cylinder radius, offsets, length, wall thickness,
and readout radius. This is useful separability and a severe physical gate:
the schedule discriminates a shell-coded observable already sitting near
background amplitude, and leaves the signal amplitude wherever the physical
readout model places it. The calibrated Stage 3 Casimir/metric readout scale
sits many orders below ordinary material and electromagnetic backgrounds.
Under the modeled nuisance conditions, that gap closes the direct measurement
route for the current apparatus. Any future positive laboratory claim must
first identify a concrete material, propagation, or readout transduction
mechanism that creates a shell-coded observable at order-unity
signal-to-background ratio and survives the same perturbation test.

\section{Interpretation}

Within these role-resolved proxy layers, the White et al. sphere-cylinder
geometry produces a near-shell Casimir boundary morphology. The mid-scale
scalar channel gives the most spatially discriminating read: signed proxy
magnitude concentrates in the annular source band, far-field magnitude is
small, boundary participation remains secondary, and the radial peak stays near
the middle of the sphere-wall gap. The tensor-scale source proxy preserves the
same source-shell placement and keeps the central readout channel small.

The broad-window field describes a separate scale channel. Larger loops
connect the apparatus over most of the sampled section, producing
apparatus-wide signed scalar participation. The broad field retains a
shell-enriched largest-magnitude component and combines localized sphere-wall
morphology with long-scale apparatus connection. The Alcubierre-like comparison
is carried by the separated mid-scale pair interaction, which localizes around
the intended shell band.

The readout-axis metrics describe overlap between the proxy fields and central
accounting masks. In the clean mid-scale channel, most signed proxy magnitude
is assigned to the near-shell band. The explicit transit/readout mask receives
about one tenth of one percent of the scalar magnitude and zero calibrated
tensor-scale readout source fraction. Larger scale channels add weak central
participation as part of the broader apparatus field. The source proxy is
concentrated where a shell would be drawn; the central path appears mainly
through apparatus-scale accounting.

The source-demand comparison supplies the scale-limiting result. The
Alcubierre comparison requires stress components, sign placement, magnitude
comparison, and metric source demand in addition to scalar morphology. In the
present calibrated proxy, the morphology is strong enough to be a real
shell-placement result, and the source magnitude sits many orders beneath a
representative Alcubierre demand. Recent work by Garattini and Zatrimaylov
frames this same layer in terms of energy conditions and metric-source
structure for warp-drive geometries \cite{GarattiniZatrimaylov2024}. A
proposed transit-time or current/photon/electron readout also carries ordinary
electromagnetic competition: capacitance, impedance, dispersion, material
response, contact loading, patch potentials, surface roughness, and
geometry-dependent leakage can all dominate a laboratory signal.

The synthetic readout-discrimination layer turns that laboratory statement into
a gate. Geometry modulation can separate a shell-coded template from nuisance
templates once the shell-coded observable reaches near-background scale; the
current block-bootstrap geometry-derived run places that recovery threshold
near signal-to-background ratio $1.2$. This result gives the perturbation
schedule a precise role. It is a discriminator across channels, with cylinder
radius, sphere size, offsets, wall thickness, length, and readout radius acting
as the modulation coordinates. The calibrated Casimir/metric readout scale from
Stage 3 sits far below this recovery threshold. Taken together, Stage 3 and
Stage 4 give a negative direct-readout result for the modeled apparatus: the
shell-shaped morphology remains real in the proxy, and the corresponding
physical readout amplitude falls short of the level where the schedule can
recover it. The perturbation schedule is therefore a diagnostic for any
proposed transduction mechanism. A viable experimental claim must first show a
concrete propagation/readout channel that produces a shell-coded observable at
order-unity signal-to-background ratio and preserves the geometry-coded
template.

\section{Code and Data Availability}

Code and generated data for this study are available in the GitHub repository
\url{https://github.com/dgoldman0/spacetime-metric-engineering}. The relevant
package is \path{white_casimir_source_function_audit} in the repository's
\path{white_casimir_intersection/experiments} tree. That package contains the
scalar fidelity case table, aggregate summary, mean field arrays, tensor-scale
harness, finite-difference stress-channel proxy, Alcubierre source-demand
comparison, linearized metric-bound estimate, ordinary-EM proxy, and pytest
checks for the local crossing geometry, scoring, bookkeeping, and Stage 3
harness.

The focused tensor-scale output used for the source-function tables is stored
under the external research run directory:
\begin{quote}
\small
\path{/media/kir/9CDCBD3EDCBD140C/Research/}\\
\path{white_casimir_tensor_scale_probe/stage3_focused_20260530}
\end{quote}
The compact output files are:
\begin{itemize}
\item \path{calibration/calibration_summary.csv}
\item \path{tensor/stress_channel_summary.csv}
\item \path{tensor/bootstrap_summary.csv}
\item \path{scale/source_demand_ratio.csv}
\item \path{scale/linearized_metric_bound.csv}
\item \path{material/material_competitor_summary.csv}
\item \path{readout/readout_coupling_summary.csv}
\end{itemize}

The synthetic readout-discrimination output used for the recovery-threshold
table is stored under:
\begin{quote}
\small
\path{/media/kir/9CDCBD3EDCBD140C/Research/}\\
\path{white_casimir_synthetic_discrimination/}\\
\path{stage4_block_bootstrap_geometry_sbr_refine_20260530}
\end{quote}
The compact and data-heavy output files are:
\begin{itemize}
\item \path{summary.json}
\item \path{synthetic/template_matrix.parquet}
\item \path{synthetic/shell_template_samples.parquet}
\item \path{synthetic/synthetic_observations.parquet}
\item \path{synthetic/recovery_ledger.parquet}
\item \path{synthetic/false_positive_ledger.parquet}
\item \path{synthetic/schedule_recommendation.csv}
\end{itemize}

\section{Conclusion}

The audit gives White et al.'s sphere-cylinder observation a narrower and
cleaner physical identity. The apparatus behaves, in these role-resolved
proxies, as a structured Casimir boundary system. The selective scalar field
and tensor-scale source proxy both place the strongest accounting in the
annular sphere-wall band, with far-field leakage small and the readout path
retained as a minor measurement channel. The Alcubierre resemblance therefore
has a concrete basis at the morphology/source-placement layer: the geometry can
shape vacuum-boundary proxy structure into the same kind of shell role drawn in
the metric comparison.

The scale and readout layers are the gate. The calibrated proxy places the
gravitational/source magnitude many orders beneath representative Alcubierre
source demand and gives a linearized timing bound of order
$10^{-76}\,\mathrm{s}$. Stage 4 then asks whether geometry-coded perturbations
could recover a shell template from ordinary electromagnetic, material, drift,
vibration, and readout backgrounds. Under the block-bootstrap,
geometry-derived nuisance model, recovery becomes reliable near
signal-to-background ratio $1.2$, with an ordinary-EM-only false-positive rate
near $0.034$. That is a severe threshold because the calibrated
Casimir/metric readout channel sits far below the electromagnetic/material
background scales. The combined Stage 3/Stage 4 result is therefore a negative
direct-detection result for the modeled apparatus.

The synthetic schedule remains useful in a narrower role. It can rank geometry
perturbations and test whether a proposed readout channel carries a
shell-coded response with enough amplitude and template orthogonality. It is a
discriminator, with signal amplitude supplied entirely by the physical
transduction model being tested. Any future positive laboratory claim must
first supply an independently modeled material, propagation, or readout
mechanism that lifts a shell-coded observable to order-unity
signal-to-background ratio and passes the same ordinary-EM-only false-positive
test.

White's geometry is therefore valuable as a reproducible shell-forming Casimir
configuration and as a falsification harness for readout proposals. Its warp
relevance, in this audit class, is geometric and organizational: it tells us
where the source-like proxy lives, how sharply it stays in the shell role, how
small the associated metric scale is after SI calibration, and where a
geometry-coded laboratory schedule begins to recover a signal. The present
claim ends at that gate: shell-shaped Casimir boundary morphology with
calibrated source scale, material competition, synthetic recovery threshold,
and a closed direct laboratory detection path under the modeled backgrounds.

\begin{thebibliography}{9}

\bibitem{White2021}
H. White, J. Vera, A. Han, A. R. Bruccoleri, and J. MacArthur,
``Worldline numerics applied to custom Casimir geometry generates
unanticipated intersection with Alcubierre warp metric,''
\emph{European Physical Journal C} \textbf{81}, 677 (2021).
\href{https://doi.org/10.1140/epjc/s10052-021-09484-z}{doi:10.1140/epjc/s10052-021-09484-z}.

\bibitem{Alcubierre1994}
M. Alcubierre,
``The warp drive: hyper-fast travel within general relativity,''
\emph{Classical and Quantum Gravity} \textbf{11}, L73--L77 (1994).
\href{https://doi.org/10.1088/0264-9381/11/5/001}{doi:10.1088/0264-9381/11/5/001}.

\bibitem{Gies2006}
H. Gies and K. Klingmuller,
``Worldline algorithms for Casimir configurations,''
\emph{Physical Review D} \textbf{74}, 045002 (2006).
\href{https://doi.org/10.1103/PhysRevD.74.045002}{doi:10.1103/PhysRevD.74.045002}.

\bibitem{GarattiniZatrimaylov2024}
R. Garattini and K. Zatrimaylov,
``Black Holes, Warp Drives, and Energy Conditions,''
\emph{Physics Letters B} \textbf{856}, 138910 (2024).
\href{https://doi.org/10.1016/j.physletb.2024.138910}{doi:10.1016/j.physletb.2024.138910}.

\bibitem{GoldmanRail2026}
D. S. Goldman,
``From the Warp/Wormhole Interface to a Throat-Supported Shift Rail:
A Packet-Centered Active-Rail Architecture after the Wormhole-Warp-Drive
Correspondence,'' preprint (2026).
\href{https://doi.org/10.13140/RG.2.2.10402.39368}{doi:10.13140/RG.2.2.10402.39368}.

\bibitem{GoldmanDisclosure}
D. S. Goldman,
``Active-Rail Spacetime Metric Engineering System: Technical Disclosure and
Source-Placement Architecture,'' technical disclosure (2026).
\href{https://github.com/dgoldman0/spacetime-metric-engineering/blob/main/active_rail_technical_disclosure.pdf}{repository PDF}.

\end{thebibliography}

\end{document}
